                      IMC 2017, Blagoevgrad, Bulgaria


                                  Day 1, August 2, 2017




Problem 1.    Determine all complex numbers λ for which there exist a positive integer n and
a real n × n matrix A such that A2 = AT and λ is an eigenvalue of A.
                                (Proposed by Alexandr Bolbot, Novosibirsk State University)
Solution.   By taking squares,
                            A4 = (A2 )2 = (AT )2 = (A2 )T = (AT )T = A,

so
                                            A4 − A = 0;
it follows that all eigenvalues of A are roots of the polynomial
                                                          √      X4 − X.
    The roots of X 4 − X = X(X 3 − 1) are 0, 1 and −1±2 3i . In order to verify that these values
are possible, consider the matrices
                                                                       0   0  0    0
                                                                                      
                                                     √   !
                                                    3
                                           −√12  2
                                                                      0   1  0    0 
                                                                                  √ 
          A0 = 0 ,      A1 = 1 ,     A2 =                    ,   A4 =         1    3.
                                             − 23 − 12                0   0 −√2   2
                                                                       0   0 − 23 − 12

The numbers√0 and 1 are the eigenvalues of the 1 × 1 matrices A0 and A1 , respectively. The
numbers −1±2 3i are the eigenvalues of A2 ; it is easy to check that
                                                     √ !
                                            −  1
                                                 − 23
                                    A22 =   √2
                                              3
                                                           = AT2 .
                                             2
                                                 − 12

The matrix A4 establishes all the four possible eigenvalues in a single matrix.


Remark. The matrix A2 represents a rotation by 2π/3.



Problem 2.    Let f : R → (0, ∞) be a dierentiable function, and suppose that there exists
a constant L > 0 such that
                                     f 0 (x) − f 0 (y) ≤ L x − y
for all x, y . Prove that                       2
                                         f 0 (x) < 2Lf (x)
holds for all x.
                                                (Proposed by Jan ’ustek, University of Ostrava)




                                                 1
Solution.    Notice that f 0 satises the Lipschitz-property, so f 0 is continuous and therefore
locally integrable.
    Consider an arbitrary x ∈ R and let d = f 0 (x). We need to prove f (x) > 2L d2
                                                                                    .
    If d = 0 then the statement is trivial.
    If d > 0 then the condition provides f 0 (x−t) ≥ d−Lt; this estimate is positive for 0 ≤ t < Ld .
By integrating over that interval,
                                                Z d                          Z d
                                                        L                         L                   d2
               f (x) > f (x) − f (x − Ld ) =                f 0 x − t dt ≥
                                                                     
                                                                                      (d − Lt)dt =       .
                                                    0                         0                       2L
     If d < 0 then apply f (x + t) ≤ d + Lt = −|d| + Lt and repeat the same argument as
                            0


                                           Z |d|                              Z |d|
                                                L                                     L                          d2
           f (x) > f (x) − f (x + |d|                         0
                                                                   
                                   L
                                      )=                − f (x + t) dt ≥                  (|d| − Lt)dt =            .
                                            0                                     0                              2L

Problem 3.     For any positive integer m, denote by P (m) the product of positive divisors of m
(e.g. P (6) = 36). For every positive integer n dene the sequence
                     a1 (n) = n,     ak+1 (n) = P (ak (n)) (k = 1, 2, . . . , 2016).
   Determine whether for every set S ⊆ {1, 2, . . . , 2017}, there exists a positive integer n such
that the following condition is satised:
       For every   k with 1 ≤ k ≤ 2017, the number ak (n) is a perfect square if and only if k ∈ S .
                                                            (Proposed by Matko Ljulj , University of Zagreb)
Solution.   We prove that the answer is yes; for every S ⊂ {1, 2, . . . , 2017} there exists a
suitable n. Specially, n can be a power of 2: n = 2w1 with some nonnegative integer w1 . Write
ak (n) = 2wk ; then
                                                                                                wk (wk +1)
               2wk+1 = ak+1 (n) = P (ak (n)) = P (2wk ) = 1 · 2 · 4 · · · 2wk = 2                    2       ,
so
                                                            wk (wk + 1)
                                           wk+1 =                       .
                                                                 2
The proof will be completed if we prove that for each choice of S there exists an initial value
w1 such that wk is even if and only if k ∈ S .
Lemma. Suppose that the sequences (b1 , b2 , . . .) and (c1 , c2 , . . .) satisfy bk+1 =              and
                                                                                           bk (bk +1)
                                                                                                2
ck+1 = ck (ck +1)
            2
                  for k ≥ 1, and c1 = b1 + 2 . Then for each k = 1, . . . m we have ck ≡ bk + 2
                                            m                                                       m−k+1

(mod 2m−k+2 ).
   As an immediate corollary, we have bk ≡ ck (mod 2) for 1 ≤ k ≤ m and bm+1 ≡ cm+1 + 1
(mod 2).
Proof. We prove the by induction. For k = 1 we have c1 = b1 + 2                so the statement holds.
                                                                            m

Suppose the statement is true for some k < m, then for k + 1 we have
                                                         
            ck (ck + 1)      bk + 2m−k+1 bk + 2m−k+1 + 1
     ck+1 =             ≡
                  2                         2
             2      m−k+2        2m−2k+2
            bk + 2        bk + 2         + bk + 2m−k+1
          =                                            =
                                   2
            bk (bk + 1)                                  bk (bk + 1)
          =             + 2m−k + 2m−k+1 bk + 22m−2k+1 ≡              + 2m−k                         (mod 2m−k+1 ),
                  2                                            2
therefore ck+1 ≡ bk+1 + 2m−(k+1)+1 (mod 2m−(k+1)+2 ).

                                                             2
    Going back to the solution of the problem, for every 1 ≤ m ≤ 2017 we construct inductively
a sequence (v1 , v2 , . . .) such that vk+1 = vk (v2k +1) , and for every 1 ≤ k ≤ m, vk is even if and
only if k ∈ S .
    For m = 1 we can choose v1 = 0 if 1 ∈ S or v1 = 1 if 1 ∈/ S . If we already have such a
sequence (v1 , v2 , . . .) for a positive integer m, we can choose either the same sequence or choose
                                                              v 0 (v 0 +1)
v10 = v1 + 2m and apply the same recurrence vk+1       0
                                                            = k 2k . By the Lemma, we have vk ≡ vk0
(mod 2) for k ≤ m, but vm+1 and vm+1 have opposite parities; hence, either the sequence (vk )
or the sequence (vk0 ) satises the condition for m + 1.
    Repeating this process for m = 1, 2, . . . , 2017, we obtain a suitable sequence (wk ).

Problem 4.   There are n people in a city, and each of them has exactly 1000 friends (friendship
is always symmetric). Prove that it is possible to select a group S of people such that at least
n/2017 persons in S have exactly two friends in S .
        (Proposed by Rooholah Majdodin and Fedor Petrov, St. Petersburg State University)
Solution.   Let d = 1000 and let 0 < p < 1. Choose the set S randomly such that each people
is selected with probability p, independently from the others.
    The probability that a certain person is selected for S and knows exactly two members of
S is                                      
                                             d 3
                                       q=      p (1 − p)d−2 .
                                             2
   Choose p = 3/(d + 1) (this is the value of p for which q is maximal); then
                                         3      d−2
                                    d    3       d−2
                                q=                        =
                                    2   d+1      d+1
                                          −(d−2)
                 27d(d − 1)           3               27d(d − 1) −3         1
               =           3
                                1+                 >            3
                                                                  ·e >        .
                  2(d + 1)          d−2                2(d + 1)         2017
   Hence, E |S| = nq > 2017
                         n
                             , so there is a choice for S when |S| > 2017
                                                                      n
                                                                          .
               


Problem 5.     Let k and n be positive integers with n ≥ k2 − 3k + 4, and let
                                f (z) = z n−1 + cn−2 z n−2 + . . . + c0

be a polynomial with complex coecients such that
                               c0 cn−2 = c1 cn−3 = . . . = cn−2 c0 = 0.

Prove that f (z) and z n − 1 have at most n − k common roots.
               (Proposed by Vsevolod Lev and Fedor Petrov, St. Petersburg State University)
Solution.  Let M = {z : z n = 1}, A = {z ∈ M : f (z) 6= 0} and A−1 = {z −1 : z ∈ A}. We have
to prove |A| ≥ k.
   Claim.

                                            A · A−1 = M.
That is, for any η ∈ M , there exist some elements a, b ∈ A such that ab−1 = η .
  Proof. As is well-known, for every integer m,


                                               n if n|m
                                              (
                                     X
                                         zm =
                                     z∈M
                                               0 otherwise.

                                                  3
Dene cn−1 = 1 and consider
           X                       X             n−1
                                                 X                n−1
                                                                  X                 n−1 X
                                                                                    X   n−1               X
                 2                           2                j              `
                z f (z)f (ηz) =          z             cj z             c` (ηz) =             cj c` η `         z j+`+2 =
          z∈M                      z∈M           j=0              `=0               j=0 `=0               z∈M

          n−1 X
          X   n−1            X n if n|j + ` + 2            n−2
                                                             X
                         `                           2
      =              cj c` η                      = cn−1 n +     cj cn−2−j η n−2−j n = n 6= 0.
          j=0 `=0
                                0    otherwise                                      j=0
                             z∈M

Therefore there exists some b ∈ M such that f (b) 6= 0 and f (ηb) 6= 0, i.e. b ∈ A, and
a = ηb ∈ A, satisfying ab−1 = η .
   By double-counting the elements of M , from the Claim we conclude
                |A| |A| − 1 ≥ M \ {1} = n − 1 ≥ k 2 − 3k + 3 > (k − 1)(k − 2)
                           


which shows |A| > k − 1.




                                                                    4
